% PDF-derived LaTeX transcription. Original manuscript source was unavailable.
% Published PDF SHA-256: 56b674defee5a0689d36b63499e7b32e285a46a510b531b15168df789ace202f
% Extracted text SHA-256: 6ba7c70fdb98290b5bc913601fdcf9e7eb7287e492868433c2d346ede725ab3b
\documentclass[11pt]{article}
\usepackage[margin=1in]{geometry}
\usepackage{amsmath,amssymb}
\setlength{\emergencystretch}{3em}
\title{Physical-response diagnostics, ensemble agreement, and fixed-coordinate fidelity in residue-level IDP simulations}
\author{}
\date{}
\begin{document}
\maketitle
\noindent\textbf{PDF-derived LaTeX transcription.} This editable text follows the PDF extraction. Refer to the published PDF for figures, table layouts, and mathematical notation.

% Extracted chunk 1

Physical-response diagnostics, ensemble agreement, and
fixed-coordinate fidelity in residue-level IDP simulations

\section*{Abstract}

Reliable use of residue-level simulations requires separate evidence for physical response,
sampled ensembles, and numerical implementation. We examine these questions through a ret-
rospective comparison of a development study of the ACTR/NCOA3 transactivation domain,
a fixed-condition four-construct CAL V ADOS2 panel, and deposited-coordinate reference diag-
nostics. Retained ACTR reports contain twelve 200 ns HPS-Urry replicas at two temperatures.
Their second-half equal-bead radius-of-gyration means increase from 2.271425 to 2.346678 nm,
opposite to the decreasing direction specified in the development provenance. Temperature is
confounded with execution node, and the study does not establish a causal temperature ef-
fect or a completed physical acceptance test. Separately, independent reductions of 240,000
saved samples from 24 CAL V ADOS2 replicas reproduce the frozen four-construct sampling and
comparison criteria after a 20\% cut. The largest rank-normalized split diagnostic is 1.0001992,
and the median absolute relative radius error against the supplied reference values is 4.7953\%.
PTMA lies outside the preferred uncertainty-ratio band while the aggregate panel still sat-
isfies its predeclared rule. A reported private 32-frame cross-engine energy replay and later
coordinate-readback diagnostics address different numerical questions; their original outputs
are not available here for independent reproduction. Published OpenABC reference tables
provide supporting methodological context, but cannot authenticate that private replay. The
combined evidence shows why agreement within one validation target cannot qualify another.
It supports bounded ensemble agreement for the specified panel and preserves a separate de-
velopment discrepancy and reference-reproduction limits, without asserting a new force field or
general physical qualification.

\section*{1 Introduction}

Intrinsically disordered protein (IDP) simulations are often evaluated through a compact set of
observables, including the radius of gyration, Rg. Agreement in such an observable can arise under
several scientifically different circumstances. Replicas may agree because they sample the same
distribution, because they retain a common initialization bias, or because the observable does
not resolve a relevant difference. A model may reproduce an average size at a fixed condition
while failing to supply the intended evidence about a change of conditions. Numerical agreement
between two implementations at fixed coordinates tests another property: whether the exercised
energy expressions and conventions agree on those inputs.
These distinctions matter when simulations will support subsequent representation learning or
scientific inference. A finite trajectory, an apparently stable mean, and an energy comparison are
useful observations, but each has a specific denominator and interpretation. Combining them into
a single notion of a successful simulation obscures which requirement was actually examined. Prac-
tical sampling guidance therefore emphasizes diverse starting conditions, time-series diagnostics,
correlation and uncertainty, while retaining the possibility of states that the observed trajectories
did not visit [ 3].
1

We study this problem using three evidence sets from an existing residue-level simulation inves-
tigation. First, an ACTR HPS-Urry development panel recorded a radius response under two nom-
inal temperatures. Second, a separate CAL V ADOS2 panel assessed four other constructs at fixed,
construct-specific conditions using frozen sampling and experimental-comparison rules. Third,
a retained research inventory described a deposited-coordinate energy replay and later numeri-
cal readback investigations. The available evidence is uneven: original reports and configuration
sources support the ACTR summaries; complete saved arrays permit independent re-reduction of
the CAL V ADOS2 panel; the private replay and later observer outcomes remain source-reported
summaries.
Our contribution is the joint analysis of these distinct validation targets and their actual ev-
idence levels. We retain both the ACTR direction discrepancy and the positive fixed-condition
panel. We also examine why a fixed-coordinate reference comparison, even if successful within its
own scope, would not resolve either ensemble question. No inherited Hamiltonian, simulation pack-
age, sampling diagnostic, or reference dataset is introduced as a new method here. In particular,
the four-construct study is neither an ACTR rerun nor a universal validation of residue-level IDP
models.

\section*{2 Related work and validation targets}

OpenABC provides an OpenMM-based implementation environment for existing coarse-grained
biomolecular models [ 1]. Its published validation includes comparisons of fixed-configuration ener-
gies across OpenMM and reference engines: HPS against HOOMD-Blue, Mpipi against LAMMPS,
and MOFF against GROMACS. The accompanying supplement specifies model conventions and
provides numerical tables [ 2]. This is an appropriate precedent for implementation comparisons,
with the important qualification that the compared configurations had themselves been generated
by earlier simulations. A replay with zero integration calls describes the reevaluation stage, not a
claim that its source coordinates were created without dynamics.
The inspected OpenABC supplement describes HPS as combining nearest-neighbor harmonic
bonds, screened electrostatics, and an Ashbaugh–Hatch contact interaction. It attributes the Urry-
scale parameter choices to earlier work. The CAL V ADOS 2.0 source archive supplies the inherited
source and parameter context for the separate fixed-condition study considered here [ 4]. We use
these sources for attribution and model context. Their published or archived results do not establish
the behavior of a different private run without a matching implementation, configuration, and
observable definition.
Uncertainty in a sampled observable is distinct from error in the interaction model. Correlated
time observations cannot be counted as independent preparations, and nominal uncertainty intervals
describe only the variation represented by their construction [ 3]. The present panel uses a particular
fixed sequence of cuts, convergence diagnostics, effective sample-size calculations, and resampling
decisions. Those numerical choices belong to its recorded protocol; they are not universal thresholds
prescribed by the cited sampling guidance.
Table 1 sets out the three questions used in this analysis. The distinction concerns both what
is measured and what evidence is available. A positive outcome in one row does not remove an
unresolved requirement in another.
2

Table 1: Validation targets and evidence available for this retrospective analysis. The three studies
use different systems, inputs, and observables.
Target Scientific question Evidence used here
Physical-response
diagnostic
Does the recorded ACTR size re-
sponse follow the direction speci-
fied for the development compari-
son?
Original development report, prove-
nance, configuration and runner;
independent arithmetic on saved
replica summaries.
Fixed-condition
ensemble agree-
ment
Do six replicas per construct sat-
isfy the specified sampling and
reference-comparison rules?
Frozen protocol, all saved scalar
and coordinate arrays, receipt
chains and independently imple-
mented reductions.
Fixed-coordinate
fidelity
Do exercised implementations
agree on deposited configurations,
and where does a later numerical
discrepancy appear?
Published OpenABC methods and
tables; private replay and observer
outcomes retained at inventory-
summary level.

\section*{3 Methods}


\section*{3.1 Retrospective evidence and scope}

No molecular dynamics, checkpoint inference, or fresh cross-engine replay was executed for this
manuscript. Independent calculations reduce existing reports and saved arrays. For ACTR, the cur-
rent analysis checks replica-summary arithmetic and source/configuration consistency; raw ACTR
coordinate traces and their segment records were not available in the current evidence collection.
For the four-construct panel, all saved chunks, associated receipt chains, and the exact original
protocol were available. Independent reductions check the scalar time series, coordinate-based Rg,
sampling statistics, and the aggregate decision.
The private deposited-coordinate replay and subsequent observer results require a different
evidentiary treatment. Their values are included as findings reported by the retained research
inventory, not as independently reproduced numerical measurements. Original unrounded outputs,
assertion records, and coordinate/parameter arrays for those particular studies were not captured.
Published OpenABC validation is considered separately, rather than substituted for those private
artifacts.
The experimental reference numbers used by the private studies are retained inputs. For ACTR,
the captured provenance reports their literature origin and classifies the simulation as development
rather than held-out validation. Exact experimental conditions, raw scattering curves, experi-
mental uncertainties, terminal chemistry, and the stated parameterization overlap have not been
independently verified from the full primary article bodies in the current evidence set. This analysis
preserves that provenance classification without presenting it as a newly authenticated experimental
benchmark.
3


\section*{3.2 Observable definitions}

For coordinates xi of N residue beads, define the equal-bead and mass-weighted centers and radii
by
x = 1
N
NX
i=1
xi, Rg,eq =
"
1
N
NX
i=1
\ensuremath{\parallel}xi \ensuremath{-} x\ensuremath{\parallel}2
2
\#1/2
, (1)
xcm =
P
i mixiP
i mi
, Rg,mw =
[U+0014] P
i mi\ensuremath{\parallel}xi \ensuremath{-} xcm\ensuremath{\parallel}2
2P
i mi
[U+0015]1/2
. (2)
The ACTR comparison uses second-half equal-bead diagnostic means. The fixed-condition panel
uses mass-weighted Rg after its specified cut. These definitions and averaging windows are kept
separate throughout. For the panel’s source-specific residue masses, the N- and C-terminal mass
increments are 2 and 16 atomic mass units, respectively. The coordinate reduction uses the saved
Cartesian arrays directly; it does not introduce a new alignment, minimum-image transformation,
or periodic rewrapping convention.
Let Rr denote the time mean retained for replica r. Within a condition with R replicas, we
report
R = 1
R
RX
r=1
Rr, s 2
rep = 1
R \ensuremath{-} 1
RX
r=1
(Rr \ensuremath{-} R)2. (3)
The ACTR standard deviation describes spread of the six recorded replica means. It is not a
confidence interval or an estimate obtained by treating individual frames as independent. The
fixed-condition panel estimates uncertainty in its mean using the explicit hierarchical resampling
procedure below.

\section*{3.3 ACTR development comparison}

The retained configuration specifies a 71-residue ACTR/NCOA3 transactivation-domain sequence,
with three compact and three extended starts at each of 278.15 and 318.15 K. Each of the twelve
replicas targets 200 ns, for 2,400 ns aggregate. A 10 fs integration step and observations every 5,000
steps yield 4,000 saved frames per replica, or 48,000 across both conditions. The reported size
comparison uses the second half of each run as a development diagnostic. This window is not
asserted to be a validated equilibration cutoff.
The configured interaction model is OpenABC HPS-Urry, with contact energy scale 0.2 kcal/mol,
\ensuremath{\mu} = 1 , \ensuremath{\Delta} = 0 .08, relative dielectric constant 80, and model ionic strength 0.2 M. The inspected
runner computes a condition-specific Debye screening length from temperature, ionic strength, and
dielectric constant, and supplies it to the screened electrostatic interaction. Thus temperature
also changes this configured screening length. The recorded integrator uses a friction of 1 ps\ensuremath{-}1,
a 40 nm periodic box, and mixed-precision CUDA execution. The source configuration includes
residue charges without explicit terminal charges, because the experimental terminal chemistry
is unresolved. These are the private run’s conventions, not a claim of exact reproduction of all
experimental conditions.
Crucially, the two temperatures are assigned to different execution nodes. The recorded design
therefore does not isolate a temperature effect from node or backend effects. The intended qualita-
tive direction is a smaller Rg at the higher temperature, as specified by the retained development
provenance. The report explicitly leaves the formal scientific direction assessment unscored. No
equilibration, effective sample size, rank-normalized convergence, or raw-scattering goodness-of-fit
result from the separate panel is transferred retrospectively to ACTR.
4


\section*{3.4 Fixed-condition CAL V ADOS2 panel}

The second study contains DSS1, PTMA, NHE6cmdd, and ANAC046, with conditions recorded in
Table 2. Each construct has six replicas: three compact and three extended starts with separate
recorded seeds. The protocol uses a Langevin middle integrator, a 10 fs step, friction 0.01 ps\ensuremath{-}1, and
double-precision CUDA execution. Each replica contains 108 steps, corresponding to 1,000 ns, and
stores observations every 10,000 steps. Thus there are 10,000 saved samples per replica, 60,000 per
construct, and 240,000 over the panel. These durations describe the original recorded study, not
new dynamics in this work.
All four constructs use the source’s single-chain cubic periodic-box rule
L = [0.38(N \ensuremath{-} 1) + 4] nm. (4)
The first 20\% of each scalar series is discarded, leaving 8,000 observations per replica and 192,000
across the panel. The corresponding retained duration is 800 ns per replica. The protocol de-
scribes the constructs as a fixed-condition held-out panel; that designation is retained as a protocol
classification and is not generalized to independence from every model-development dataset.
Table 2: Construct-specific inputs in the frozen CAL V ADOS2 protocol. I is ionic strength. The
reference Rg and uref are supplied experimental values and uncertainties, not independently remea-
sured experimental data.
Construct Residues T (K) pH I (M) Rg,ref \ensuremath{\pm} uref (nm)
DSS1 71 288 7.4 0.17 2.5 \ensuremath{\pm} 0.1
PTMA 111 288 7.4 0.16 3.7 \ensuremath{\pm} 0.2
NHE6cmdd 116 288 7.4 0.17 3.2 \ensuremath{\pm} 0.2
ANAC046 167 298 7.0 0.14 3.6 \ensuremath{\pm} 0.3

\section*{3.5 Frozen sampling and uncertainty rules}

For each post-cut replica, the drift diagnostic compares the first and last of four equal chronological
blocks. Their absolute mean difference, divided by the absolute full retained mean, must not exceed
0.10. A separate split-half comparison, normalized by that retained mean, must not exceed 0.05 in
any replica. Both normalizations use an absolute-mean floor of 10\ensuremath{-}12 nm in the recorded assessor.
The rank-normalization and folding approach follows Vehtari et al. [ 5]; the thresholds and
single-replica ESS rule below remain choices of the frozen study protocol. The rank-normalized
diagnostic splits each of the six retained series into two equal halves. If rj is the average rank of
a value among the pooled split observations and M is the pooled count, the assessor maps ranks
through \ensuremath{\Phi}\ensuremath{-}1[(rj \ensuremath{-} 3/8)/(M + 1/4)], where \ensuremath{\Phi} is the standard-normal cumulative distribution. The
diagnostic is computed both on these normalized values and on ranks of absolute deviations from
the pooled median. For either normalized split matrix, with m chains of length n, the conventional
variance ratio used by the frozen assessor is
W = 1
m
mX
r=1
s2
r, B = n Varr(zr), bR =
[U+0014] (n \ensuremath{-} 1)W/n + B/n
W
[U+0015]1/2
. (5)
Here the within-chain variances and the variance of chain means use sample denominators. The
larger bulk or folded value must be at most 1.05. In this study m = 12 and n = 4 , 000 after
5

splitting. These diagnostics test the observed scalar series under the fixed rule; passing them does
not prove complete exploration of configuration space.
For a post-cut series of length n0 = 8 , 000, the effective sample-size (ESS) calculation uses
centered lag correlations. The original assessor obtains autocovariances by a fast Fourier transform
and divides each lag product sum by the number of available pairs. Consecutive positive lag pairs
are accumulated until the first nonpositive pair. Writing K for the last included pair, its rule is
\ensuremath{\tau} = 1 + 2
KX
k=1
(b\ensuremath{\rho}2k\ensuremath{-}1 + b\ensuremath{\rho}2k), N eff = max\{1, min(n0, n0/\ensuremath{\tau} )\}. (6)
The minimum is 200 per replica and 1,200 after summing the six replica estimates. The independent
reduction uses direct lag-dot products for the same definition, rather than importing the original
analysis source. Clipping at the observed count is part of this particular estimator.
Start-state agreement first requires a relative separation of at least 0.30 between the recorded
compact and extended initial Rg group means. This initial separation is their extended-minus-
compact mean difference divided by the average of those two means, with the same positive nu-
merical floor. For the post-cut means, let Rc and Re be the averages over the three compact and
three extended starts. Their absolute difference must not exceed
A = max
(
2
r
s2c
3 + s2e
3 , 0.05

% Extracted chunk 5


Rc + Re
2

% Extracted chunk 9


)
, (7)
where s2
c and s2
e describe variation among the respective replica means. Initial size values are
retained progress records; compact initial configurations were not regenerated here.
For uncertainty in the panel mean, the frozen procedure performs 2,000 seeded hierarchical
circular-block bootstrap draws. Each draw selects six replicas with replacement. For each selected
replica, circular temporal blocks are sampled with replacement, concatenated, and truncated to
its retained sample count; their mean supplies that selected replica’s contribution. The six contri-
butions are then averaged. Replica-specific block lengths are \ensuremath{\lceil}n0/Neff \ensuremath{\rceil}, restricted to [1, n0]. The
standard deviation of bootstrap means, using a sample denominator, supplies the model standard
error, smodel. The 2.5th and 97.5th percentiles give the stated interval. The independent reduction
retains the original seeds and circular-block law, with precomputed circular window sums; it does
not execute the supplied analysis code.
The comparison ratio for construct a is the frozen operational quantity
za = |Ra \ensuremath{-} Ra,ref |q
s2
a,model + u2
a,ref
, e a = |Ra \ensuremath{-} Ra,ref |
Ra,ref
. (8)
The denominator combines the bootstrap model standard error with the reported experimental
uncertainty as the protocol specifies. It does not add force-field or other systematic errors. No
calibrated hypothesis-test probability or full-scattering agreement is inferred from za. The aggregate
rule requires all four sampling outcomes, at least three za \ensuremath{\leq} 2, none above 3, median ea \ensuremath{\leq} 0.10,
and maximum ea \ensuremath{\leq} 0.20.

\section*{3.6 Independent saved-array reduction and precision}

The panel’s saved scalar arrays retain double-precision Rg, whereas the coordinate arrays are single
precision. Coordinate-based Rg was independently recomputed with the source-specific masses.
6

Agreement is evaluated under a conservative propagation bound for coordinate rounding, rather
than demanding bit identity with an observable originally computed before coordinate conversion.
The largest absolute difference in the panel is 3.210902 \ensuremath{\times} 10\ensuremath{-}6 nm; every checked coordinate/scalar
pair remains within its recorded conservative bound.
Independent statistical reductions reproduce the original sampling decisions and comparison
numbers at their stated numerical tolerances. A standard-library inverse-normal quantile is used
instead of the source’s approximation for rank normalization. Tiny resulting numerical differ-
ences are not interpreted as a different scientific outcome. These calculations establish consistency
of the retained samples and methods. They do not reproduce the original OpenMM dynamics,
independently establish every runtime parameter consumed by the engine, or qualify a different
implementation.

\section*{4 Results}


\section*{4.1 ACTR: a retained development discrepancy}

The six ACTR replica means at each condition give the summaries in Table 3. The high-minus-
low difference is +0.0752539632 nm. The experimental values reported by the retained provenance
instead decrease by 0.24 nm. Thus the observed development diagnostic has the opposite sign to
the comparison direction specified in that provenance. It does not supply the intended positive
evidence about physical response.
Table 3: Retained ACTR second-half equal-bead diagnostic summaries. Each condition has three
compact and three extended starts. The standard deviation is across six replica means. Reference
numbers are copied from development provenance, not independently checked experimental mea-
surements.
T (K) Replicas Mean Rg (nm) Replica SD (nm) Reported reference Rg (nm)
278.15 6 2.271425 0.078096 2.63
318.15 6 2.346678 0.065386 2.39
At 278.15 K, compact and extended mean sizes are 2.270728 and 2.272122 nm, respectively. At
318.15 K, the corresponding means are 2.348007 and 2.345350 nm. That proximity does not convert
the development study into an equilibrated or held-out benchmark. In particular, its reported
difference in direction coexists with unresolved experimental inputs and the temperature–node
confounding. The formal scientific assessment remains unscored, rather than being retroactively
relabeled a completed acceptance failure or success.

\section*{4.2 Four-construct panel: agreement under the original fixed rule}

All 24 CAL V ADOS2 replica series satisfy the original drift, split-half, rank-normalized, ESS, and
start-state checks. The independently reduced means and intervals appear in Table 4. The aggre-
gate median absolute relative error is 0.0479526813, and the largest relative error is 0.1226623393
for PTMA. Three constructs have z \ensuremath{\leq} 2; PTMA has z = 2 .268158, below the absolute limit of 3.
The panel therefore satisfies its frozen aggregate rule without requiring that every construct meet
the preferred band.
The model bootstrap standard errors are 0.003406, 0.006218, 0.003951, and 0.007623 nm for
DSS1, PTMA, NHE6cmdd, and ANAC046, respectively. The central estimates in Table 4 are av-
7

Table 4: Independently re-reduced mass-weighted Rg results after the specified cut. Intervals are
the frozen hierarchical-bootstrap percentile intervals for the model mean. They omit force-field and
experimental systematic uncertainty. The z ratio uses the supplied reference uncertainty through
Eq. 8.
Construct Mean Rg (nm) Model 95\% interval (nm) z 100e (\%)
DSS1 2.637509 [2.630920, 2.644308] 1.374289 5.500345
PTMA 4.153851 [4.141701, 4.165956] 2.268158 12.266234
NHE6cmdd 3.069114 [3.061511, 3.077141] 0.654303 4.090191
ANAC046 3.680021 [3.665165, 3.694945] 0.266650 2.222797
erages of the six retained replica means; bootstrap means are used to estimate uncertainty rather
than replace these central estimates. The intervals quantify uncertainty under the specified resam-
pling procedure and observed state coverage. They are not population-wide evidence of force-field
accuracy.
Table 5 shows the sampling diagnostics. The largest rank-normalized split value is approx-
imately 1.0001992. The smallest individual ESS is 2,807.45, well above the stated per-replica
threshold. The largest observed first/last-block drift is 2.6198\%, and the largest split-half differ-
ence is 1.4692\%. These are outcomes for mass-weighted Rg in these four systems and observation
windows. They do not establish convergence of every possible structural observable.
Table 5: Fixed-condition sampling diagnostics independently reduced from the saved scalar series.
Maxima are over the six replicas for each construct; the reported rank statistic is the larger bulk
or folded value. Pooled ESS is the sum of the six estimator values.
Construct bRmax Minimum ESS Pooled ESS Drift max (\%) Split-half max (\%)
DSS1 1.000071 7263.79 46645.77 0.8875 0.6005
PTMA 1.000124 4511.07 27987.91 0.7664 1.0177
NHE6cmdd 1.000063 4225.94 26972.20 1.2777 1.4692
ANAC046 1.000199 2807.45 17372.91 2.6198 1.4456
In construct order, the compact–extended retained-mean differences are 0.004753, 0.000229,
0.003451, and 0.000552 nm. They lie below the source-specific allowances of 0.131875, 0.207693,
0.153456, and 0.184001 nm. The recorded initial group relative separations are all above 1.36, rather
than merely meeting the 0.30 requirement. Initial state separation therefore supplies a stronger
check than comparing starts that were already similar in Rg, while still leaving unvisited-state and
observable-specific sampling limits.

\section*{4.3 Deposited-coordinate references and numerical observations}

The retained inventory reports a private replay of 32 deposited configurations across HPS, Mpipi,
and MOFF, with no integration or minimization calls during reevaluation. The inventory reports
that all 32 passed their upstream assertions. Its reported counts and rounded energy differences
are retained in Table 6. The original private coordinate arrays, unrounded energy terms, asser-
tion records, and complete replay implementation were unavailable for independent checking. The
table consequently reports an inventory finding; it is not a newly reproduced cross-engine bench-
mark. The inventory describes a CPU, software-version-pinned replay rather than a complete
8

executable/native-library lock or GPU qualification.
Table 6: Private deposited-coordinate replay as reported by the retained research inventory. The
original private replay outputs are not independently reproduced here. Rounded CSV maxima and
reported assertion tolerances are distinct quantities, both in kJ/mol.
Model Reference engine Frames Tolerance Rounded term maximum
HPS HOOMD-Blue 10 0.01 0.0055103629
Mpipi LAMMPS 11 0.01 0.0100
MOFF GROMACS 11 0.02 0.0198
The published OpenABC supplement separately contains ten configurations in each of its
MOFF, HPS, and Mpipi energy-comparison tables, for 30 published rows [ 2]. This published
comparison is distinct from the private inventory’s 10/11/11-frame study. The published rounded
entries also warrant care: the HPS Urry contact column contains a printed difference of 0.01 kJ/mol
in one configuration, despite a prose description of exact agreement. A displayed rounded differ-
ence cannot recover an unrounded private assertion result. Published methodological agreement
therefore supplies context without authenticating the private rerun, its software environment, or
its numerical tolerance.
The same inventory reports later zero-step coordinate observations. Among 100 position rows,
88 changed relative to commanded coordinates before first force capture, with a maximum differ-
ence of 7.105427357601002 \ensuremath{\times} 10\ensuremath{-}15 nm; first force capture introduced no additional changed rows.
A tested box-add/subtract arithmetic hypothesis matched those observed rows, but did not authen-
ticate the engine’s native ordering. Original stage-specific coordinate arrays and observer outputs
were not supplied here, so these remain reported diagnostics.
A separate full-chain force triage in the inventory records a residual under serialized binary64
parameters of
1.0375627494303984 \ensuremath{\times} 10\ensuremath{-}4 kJ/(mol nm).
Under a CUDA-storage parameter hypothesis, the reported residual is
7.822559169383161 \ensuremath{\times} 10\ensuremath{-}10 kJ/(mol nm).
That contrast narrows a possible representation explanation without proving actual runtime pa-
rameter consumption. It also lacks an authoritative matching full-HPS force-accuracy threshold
in the available evidence. A two-particle harmonic tolerance, a per-particle norm bound, or an
energy-ledger closure cannot supply that missing criterion.

\section*{5 Discussion}

The principal empirical contrast is two-sided. The ACTR report records a size response that
does not support its intended development comparison, while the independent fixed-condition re-
reduction supports the positive panel decision under its own original rules. Neither outcome removes
the other. They concern different proteins, models, conditions, definitions of Rg, averaging windows,
and evidence levels. The different ACTR node assignments further prevent a causal interpretation
based on temperature alone.
The fixed-condition result is useful because it retains both individual and aggregate outcomes.
PTMA does not meet the preferred z \ensuremath{\leq} 2 band, but the predeclared rule explicitly permits one
9

such case if none exceeds 3 and the relative-error limits hold. Reporting only an aggregate pass
would hide this structure; declaring the panel failed solely because PTMA misses the preferred band
would instead change the rule after observing the results. The appropriate statement is agreement
under the specified four-construct criterion, with the individual discrepancy visible.
Sampling uncertainty is not an uncertainty budget for the physical model. The hierarchical
bootstrap addresses variation represented by the six observed replica series and their circular tem-
poral blocks. It does not include errors in the force field, source experimental uncertainties beyond
the supplied values, experimental conditions, terminal chemistry, or states absent from the trajec-
tories. The operational combined ratio is therefore informative within the comparison contract but
is not a calibrated probability that the model is physically correct. Likewise, low scalar convergence
statistics and close starting-class means do not prove equilibration of every coordinate-dependent
observable.
Fixed-coordinate fidelity supplies an orthogonal kind of evidence. Agreement of exercised en-
ergies can constrain implementation differences at the compared coordinates and conventions. It
cannot establish a correct distribution of coordinates, a physical response to changed conditions, or
an acceptable force error when the matching reference criterion is unknown. The inherited model
and software contexts are essential to such an interpretation: coefficients, cutoffs, shifts, exclusions,
precision, and parameter representations can each change what is being compared. This paper’s
ensemble outcomes do not requalify separate contact-term or numerical-precision studies as new
independent experiments.
The tiny reported coordinate-readback differences illustrate a related distinction between locat-
ing a discrepancy and identifying its mechanism. A stage at which rows first differ is an observation
about the exercised measurement path. Matching an arithmetic hypothesis does not identify all
internal ordering or parameter-consumption operations. Numerical smallness also does not supply
a physical acceptance threshold. The original arrays and intervention records would be needed to
reproduce those reported observations and assess stronger mechanism claims.
The current evidence supports a more substantial reproduction statement for the CAL V ADOS2
panel than for the other studies. All saved samples were available and their scalar statistics and
coordinate-based radii were independently reduced. This checks retained-data consistency and the
specified decision. It remains distinct from rerunning the dynamics with independently recon-
structed software and reference inputs. For ACTR, only report-level summaries and methods were
checked. For the private replay and observers, the original inventory remains the evidence source.
These differences are retained rather than averaged into a single validation score.
Several experimental and reuse questions remain. The detailed primary ACTR experimental
and parameterization article bodies were unavailable in the captured citation set, so the specific
experimental numbers and development classification are attributed to private provenance. Raw
scattering curves, uncertainties, terminal chemistry and applicable reuse permissions remain unre-
solved. The fixed-condition panel’s reference numbers and uncertainty entries likewise retain their
supplied-source status. None of these limitations can be eliminated by a narrower sampling interval
or by agreement between two numerical implementations.

\section*{6 Conclusion}

Physical-response diagnostics, ensemble agreement, and fixed-coordinate implementation fidelity
answer different questions in residue-level IDP simulation. The retained ACTR development com-
parison gives a high-minus-low size difference of +0.075254 nm, opposite to the direction specified
by its provenance; its node–temperature confounding and unassessed convergence remain explicit.
10

A separate four-construct CAL V ADOS2 panel, independently re-reduced from all saved samples,
satisfies its original sampling and aggregate reference-comparison rules, with median relative ra-
dius error 4.7953\% and PTMA outside the preferred uncertainty-ratio band. These findings coexist
without contradiction.
The private deposited-coordinate replay and later readback/force diagnostics retain a distinct
source-reported status until their original artifacts can be checked. Published OpenABC compar-
isons support the rationale for that validation target, but do not authenticate the private results.
The resulting scientific lesson is to preserve the observable, population, protocol and evidence level
attached to each conclusion. The bounded positive panel result does not establish a new force field,
universal IDP accuracy, or a completed physical qualification of all related simulations.

\section*{7 Data and code availability}

The current study uses privately retained research reports, provenance and configuration records,
inspected method sources, saved fixed-condition scalar and coordinate arrays, and independently
authored reduction code. The original ACTR raw traces and the private deposited-coordinate
replay and observer arrays are not included in the frozen evidence supplied to this draft. Public
literature and the archived CAL V ADOS source provide their separately stated background and
method context. No public release of the private research collection, external access procedure, or
redistribution permission is established here. Supplementary record S1 retains the private source
evidence and independently authored reduction code. A separate verification report records the
remaining specific reproduction, experimental-source and access checks.

\section*{References}

[1] S. Liu, C. Wang, A. P. Latham, X. Ding, and B. Zhang. OpenABC enables flexible, simplified,
and efficient GPU accelerated simulations of biomolecular condensates. PLoS Computational
Biology 19(9), e1011442 (2023). https://doi.org/10.1371/journal.pcbi.1011442.
[2] S. Liu, C. Wang, A. P. Latham, X. Ding, and B. Zhang. Supplementary text to OpenABC en-
ables flexible, simplified, and efficient GPU accelerated simulations of biomolecular condensates
(2023). Publisher supplementary attachment to https://doi.org/10.1371/journal.pcbi.
1011442; Tables J, L and M.
[3] A. Grossfield, P. N. Patrone, D. R. Roe, A. J. Schultz, D. W. Siderius, and D. M. Zuckerman.
Best Practices for Quantification of Uncertainty and Sampling Quality in Molecular Simulations
[Article v1.0]. Living Journal of Computational Molecular Science 1(1), 5067 (2019). https:
//doi.org/10.33011/livecoms.1.1.5067.
[4] G. Tesei, T. K. Schulze, R. Crehuet, and K. Lindorff-Larsen. CAL V ADOS: Coarse-graining Ap-
proach to Liquid-liquid phase separation Via an Automated Data-driven Optimisation Scheme.
Source archive, version 2.0 (2024). https://doi.org/10.5281/zenodo.13754000.
[5] A. Vehtari, A. Gelman, D. Simpson, B. Carpenter, and P.-C. Bürkner. Rank-normalization, fold-
ing, and localization: An improved bR for assessing convergence of MCMC. arXiv:1903.08008v5
(2021). https://arxiv.org/abs/1903.08008v5. Related journal DOI: https://doi.org/10.
1214/20-BA1221.
11

\end{document}
